\subsection{Sequential Reinforcement Learning for Continuous Action Spaces}\label{seq_snn}
A key consideration when training stateful neural networks with RL, is gathering sufficiently long sequences to enable efficient training, while also allowing the hidden states to stabilize for several timesteps after initialization before calculating the gradient, which we call the warm-up period. When controlling a drone, subpar controller performance leads to too short interactions, making it challenging to achieve a baseline performance.

% \subsubsection{Bridging the Warm-Up Period}
To address this challenge in sequential SNN training, we adapt the Jump-Start Reinforcement Learning framework \cite{Uchendu2022}. They leverage a pre-trained guide policy to create a curriculum of starting conditions for a secondary policy. We implement a privileged, non-spiking actor, trained through TD3 \cite{fujimoto2018addressing}, whose primary function is to bridge the critical warm-up period required by the SNN. Both this guiding policy and the non-spiking critic receive privileged information in the form of action and observation histories. \\
% Although the actor exhibits temporal dependency, we demonstrate through importance sampling ratio analysis that this dependency does not alter the fundamental policy gradient derivation (see \autoref{sup}). \\


% \subsubsection{Temporally Training the Actor}
A replay buffer is continuously populated throughout training with transitions generated by both the guiding actor (during the initial warm-up phase) and the spiking policy (during subsequent interactions). This hybrid approach bridges the early training period, where sequences are short and therefore supervised learning from the guiding policy is effective, and the later phase, where online RL is better suited. Inspired by TD3BC \cite{FujimotoMinimal}, we incorporate a behavioral cloning (BC) term into the RL objective function that utilizes the guide policy demonstrations stored in the replay buffer. The weight of the BC term \( \lambda \) gradually decays during training, initially prioritizing the imitation of the demonstration data, progressively shifting towards reward-based optimization. The BC term in our approach serves two purposes: it avoids large changes in policy behavior, improving training stability, and it leverages the demonstration data efficiently.

The function to be maximized can be defined as:
\begin{align}
    L_{\pi} &= \mathbb{E}_{\mathbf{\tau} \sim \mathcal{D}} \left[ \sum^{100}_{i=0} Q_{\phi_1}\left( \mathbf{s}_{\tau,i}, \pi_{\theta}(\mathbf{s}_{\tau,i}| (\mathbf{s}_{\tau,i-1}, \dots, \mathbf{s}_{\tau,i})) \right) \cdot \mathbf{1}_\mathrm{i\geq50} \right] \nonumber \\
    &\quad + \lambda \mathbb{E}_{\mathbf{\tau} \sim \mathcal{D}} \left[ \sum^{100}_{i=0} \|\pi_{\theta}(\mathbf{s}_{\tau,i}| (\mathbf{s}_{\tau,i-1}...\mathbf{s}_{\tau,i})) - \mathbf{a}_{\tau,i} \|^2 \cdot \mathbf{1}_\mathrm{i\geq50} \right]
\end{align}
% function needs to be checked
\( Q_{\phi_1} \) in the first term represents the first critic network, as used in TD3. The variable \( \tau \) denotes a sequence sampled from the replay buffer, with a length of $100$. Within this sequence, \( \mathbf{s}_{\tau,i} \) and \( \mathbf{a}_{\tau,i} \) correspond to the \( i^\text{th} \) observation and action, sampled from the replay buffer, respectively. The hyperparameter \( \lambda \) controls the strength of the BC regularization and decays over time; when the guiding policy is of high quality, a larger \( \lambda \) is preferred. Finally, \( \mathbf{1}_\mathrm{i\geq50} \) equals zero during the warm-up period of 50 steps and switches to one afterward. We zero-initialize hidden states, which causes only marginal performance differences \cite{Kapturowski2019}. The critic is not time variant and thus does not need a warm-up period.

% Note that in practice, in TD3 or SAC the expected returns in the replay buffer are usually updated using newer iterations of the policy. We sample it from the transition, as recomputing the actions again requires a warm-up period for the spiking actor. While this will bias our critic towards learning the expected return from the actor which gathered the transition, it is assumed that newer policies do not diverge significantly from the transitions in the replay buffer, thanks to the BC term.

\subsection{Control of a Micro Aerial Vehicles}
We implement and evaluate our approach on a quadrotor position control task using the CrazyFlie 2.1 platform \cite{crazyflie}. The control architecture employs a spiking neural network that processes a state vector comprising position $(x, y, z)$, linear velocity $(v_x, v_y, v_z)$, orientation angles $(\theta, \phi, \psi)$, and angular velocities $(p, q, r)$. Based on these inputs, the SNN outputs a motor command for each motor, $(m_1, m_2, m_3, m_4)$ at a control frequency of $100Hz$.

SNNs work in a spike-based domain. To encode continuous values to spikes and vice versa, population based encoding and decoding are used. A linear layer encodes input values into currents, and another decodes output spikes into motor commands. \autoref{fig:task_descript} illustrates the information flow between the CrazyFlie platform and the spiking controller.

\begin{figure*}[ht!]
    \centering
    % \begin{subfigure}[b]{0.49\textwidth}
        \centering
        \includesvg[width=.49\textwidth]{figures/task_drawing.svg}
        \caption{The spiking controller receives $(x, y, z)$, linear velocity $(v_x, v_y, v_z)$, orientation angles $(\theta, \phi, \psi)$, and angular velocities $(p, q, r)$ inputs and outputs motor commands $(m_1, m_2, m_3, m_4)$. }
        \label{fig:task_descript}
\end{figure*}

The environment has a limit of $500$ timesteps before it is terminated. The guide controller is used for the first $500-N$ timesteps, after which the spiking policy interacts for the remaining $N$ steps, gradually increasing $N$ until the guide policy is only used for the warm-up period. Initially, the buffer is filled mostly with experience from the guiding policy. One runs the risk of training a spiking policy which closely resembles the guide policy, which can be undesirable when the guide is not an expert policy. Therefore, the choice of the BC term weight is crucial to prevent the spiking policy from overfitting to the guide policy.

We learn to fly in a simulated environment, using a dynamics model which has demonstrated sim-to-real bridging capabilities \cite{Eschmann2023}. The actor is penalized for errors in position, linear velocity, orientation and magnitude of actions. 
% A more detailed description of the environment and reward structure can be found in \autoref{sup:sim}.


\subsection{Spiking Actor-Critic}\label{subsection:S-AC}
Spiking neural networks leverage bio-inspired neuron models, which display a spiking behavior. The membrane potential is charged by an input current. When the membrane potential exceeds a threshold, the neuron spikes and the membrane potential is reset. This spiking behavior is non-differentiable and calls for surrogate gradients \cite{NeftciSurrogate} to enable gradient based optimization methods. We replace the threshold by a fast sigmoid, parametrized by a slope $k$, where the larger $k$, the steeper the slope. The gradient can be calculated as \autoref{eq:fast_sigmoid}:

\begin{equation}\label{eq:fast_sigmoid}
y' = \frac{1}{(1+k\cdot |x|)^2}
\end{equation}
The spiking policy employs two hidden layers of Leaky Integrate-and-Fire (LIF) neurons. Each neuron receives an input current \( I_{in} \), which incrementally charges its membrane potential \( U \). Over time, the potential decays at a rate determined by the leak factor \( \beta \). When the membrane potential exceeds a defined threshold \( U_{thr} \), the neuron emits a spike, \( s\), and its potential is subsequently reset. This spiking mechanism is described by:

\begin{equation}
U[t+1] = \beta U[t] + I_{in}[t+1] - s \cdot U_{thr}, \quad \text{where } s = \begin{cases}
1, & \text{if } U[t] > U_{thr} \\
0, & \text{otherwise}
\end{cases}
\end{equation}
% \KVDB{Its not really the equation I implement cause I first calculate U t+1 and then do the reset, but in one line it takes less space, is this fine?}


% Note that this is a first order model, which means that a large input essentially immediately causes a spike to be generated. In situations where an incoming spike should affect a neuron at a later time step, second order models could be more effective. 
% \KVDB{I DON'T KNOW IF THIS SHOULD OR SHOULD NOT BE MENTIONED, AT CONFERENCES, MANY 'OLD SCHOOL' NEUROMORPHS GOT A BIT ANNOYED BY FIRST ORDER MODELS RECEIVING SO MUCH PRAISE, AND WHEN USING EG IMU, WHERE INTEGRATION ETC IS NEEDED, SECOND ORDER MODELS ARE PROBABLY MORE EFFECTIVE? (NEXT TO FIXING SOME CONSTANTS TO 1 ;), CAN ALSO INCLUDE A FIGURE OR APPENDIX ON THIS, I HAVE SOME FIGURES ANYWAYS}

While the actor is spiking,  the critic is an ANN, which receives the aforementioned states and an action history of 32 timesteps. This asymmetric setup has been demonstrated to successfully leverage the improved training stability of ANN \cite{Tang2020, Tang2020_Hardware}.

